\scanpage{18}

On désigne par $\overline P_i$\inserted{$\subset Q_i$} l'image inverse \inserted{dans $Q_i$} du sous-grpe de torsion de $Q_i/B_i\simeq H_{i-1}$, $N_{i-1}=Q_i/\overline P_i$ est de t.f. sans torsion, donc \margin{$H_{i-1}/\mathrm{torsion}$ $\nearrow$} projectif de t.f., donc $Q_i\simeq P_i\oplus N_{i-1}$ et le complexe $K_i$ est isom à la somme directe des complexes

\diagram{\begin{tikzcd}[ampersand replacement=\&,column sep=small,row sep=small]
\text{deg} \& \& \& \text{deg }i \& i-1 \& \& \\
\overline K_i \& \cdots\arrow[r] \& 0\arrow[r] \& P_i\arrow[r,hook] \& \overline P_i\arrow[r] \& 0\arrow[r] \& 0\\
N_i[-i] \& \& 0\arrow[r] \& 0\arrow[r,dash] \& N_{i-1}\arrow[r] \& 0 \&
\end{tikzcd}}

Enfin, $P_i\subset\overline P_i$ est une inclusion de modules projectifs de t.f. tels que $\overline P_i/P_i$ $(\simeq\mathrm{tors}\ H_{i-1})$ soit de torsion, et est justiciable de la théorie des diviseurs élémentaires : \deleted{On aura} \inserted{(si $k$ \uline{principal})} ce \deleted{$\unsure{\overline K_i\simeq|K_i\simeq}$} complexe est une somme directe de complexes de la forme

\diagram{\begin{tikzcd}[ampersand replacement=\&,column sep=small,row sep=0pt]
0\arrow[r,dash] \& \overset{i}{k}\arrow[r,"n_{i,\alpha}"] \& \overset{i-1}{k}\arrow[r] \& 0 \& 1\leq\alpha\leq N_i \\
\& \& \& \& \operatorname{rg}P_i=\operatorname{rg}\overline P_i \arrow[u,equal]
\end{tikzcd}}

avec les $n_{i\alpha}\in k-(0)$ ($n_{i\alpha}$ pouvant être une unité -- donc $1$ alors) ; le nombre des $n_{i\alpha}\in k^{*}$ qui intervient est noté $\mu_i$, (il lui correspond le facteur $(\Lambda^{i}\Phi)^{\otimes\mu_i}$) ; et si on prend les \inserted{\unsure{not.|not}} $n_{\unsure{j,a|i,\alpha}}$ tels qu'ils sont croissants pour la divisibilité, on trouve une suite unique ; les multiplicités \inserted{d'occurrence} des diviseurs \unsure{él.|é}

% archive sheet number: 17
